\subsection{I.3 Improving on the Baseline}

\paragraph{Motivation.}
\textbf{Modification 1 (group size).} The baseline uses $K=2$ with the standardised advantage
$\hat A_i = (r_i - \bar r)/\sigma_r$ (see I.2). With two samples $\sigma_r$ is poorly estimated, and
a group is \emph{degenerate} whenever both completions earn equal reward ($\sigma_r=0$, advantage
ill-defined, no gradient); even non-degenerate groups give a tiny effective sample size. Raising
$K$ above $2$ is the most direct theory-motivated fix: it should reduce degenerate groups, condition
$\sigma_r$, and lower advantage variance. A team sweep of $K\in\{2,4,8,16\}$
(Table~\ref{tab:ksweep}), alongside further single-knob parameter runs (Table~\ref{tab:other-runs}),
settled the operating point at $K=8$, which gave the clearest eval gain and the most stable training.
\textbf{Modification 2 (reward).} Once $K=8$ trains stably it exposes a \emph{reward-hacking}
failure: with format and correctness weighted roughly equally, the policy formats almost perfectly
while answering only $\approx$half correctly. We downweight the format terms and replace the coarse
step-function correctness term with a denser, accuracy-weighted clipped-exponential reward, then
resume the $K=8$ checkpoint for a further $\sim$2.5k steps (5864 total). See Appendix~\ref{app:reward} for the pseudocode implementation.

\paragraph{Controlled-comparison design.}
\noindent
The comparison holds the GSM8K test split, greedy decoding, and seed-42 ordering fixed, and
evaluates the full 1319-example test set for the base model, the $K=2$ baseline, the matched
$K=8$ run, and the $K=8$ new-reward run. All three GRPO runs use the same data and the same
optimisation schedule (nominal 6516 steps, $\times 0.9 = 5864$ training steps). We normalise by
training step, not wall-clock or FLOPs: this fixes the gradient updates, schedule, and data
ordering, so any difference traces to the per-step advantage estimate. The $K=2$ baseline is the I.1
configuration trained to this same schedule. Uncertainty is a non-parametric bootstrap over the 1319
prompts (10{,}000 resamples, seed 42; per-question results in \path{evaluation/}), \emph{paired} for
cross-run differences.


\paragraph{Training dynamics and diagnostic.}
Figure~\ref{fig:dynamics} overlays the held-out curves (every 64 steps). Reward is comparable only
\emph{within} a fixed reward function, so the left panel shows the two default-reward runs and the
new-reward run is judged by KL and accuracy. The $K=8$ reward is sustained around $7.5$, whereas
$K=2$ peaks early (step $\sim$300) and collapses by step $\sim$1000 to the $-2.5$ format-penalty
floor. KL (right panel) tells the same story: $K=2$ spikes at the collapse (raw KL $\approx10$, 64-step
mean $\approx1$) and stays erratic, while both $K=8$ runs hold a bounded $\approx0.2$--$0.3$. The diagnostic
(Figure~\ref{fig:length-diagnostic}), mean eval completion length (a known GRPO failure mode)
pins the mechanism: at step $\sim$1000 $K=2$ collapses to near-empty completions, with an unstable
blow-up to $\sim$600 tokens over steps 2300--3100 before settling at a $\sim$40-token mode; both
$K=8$ runs hold a stable $\sim$300--360 tokens (new-reward $\approx362$ vs default $\approx352$).

\paragraph{Why response length, not advantage variance.}
The suggested advantage-variance diagnostic is uninformative here: $\hat A_i=(r_i-\bar r)/\sigma_r$ is
standardised \emph{per group} (Tunix sets $\hat A_i=0$ when $\sigma_r=0$), so $\mathrm{Var}(\hat A_i)\approx1$
by construction. The genuinely informative quantity, the share of degenerate groups tied to
$K_{\mathrm{eff}}$ (I.4), was not logged, so we use mean completion length, the strongest logged
GRPO-failure-mode signal.

\begin{figure}[htbp]
\centering
\optionalfigure[\linewidth]{../report-materials/data-and-figs/comparison_reward_kl.pdf}{reward and KL overlays}
\caption{Training dynamics (held-out, every 64 steps). Left: eval reward $\bar r$ for the two
default-reward runs (reward is comparable only within a fixed reward function). Right: reference-KL
$\mathrm{KL}(\pi_\theta\|\pi_{\mathrm{ref}})$ for all three runs (64-step moving average, symlog
$y$-axis).}
\label{fig:dynamics}
\end{figure}

\begin{figure}[htbp]
\centering
\optionalfigure[0.55\linewidth]{../report-materials/data-and-figs/diagnostic_response_length.pdf}{mean completion length}
\caption{GRPO-failure-mode diagnostic: mean held-out eval completion length (tokens) over training.}
\label{fig:length-diagnostic}
\end{figure}

\paragraph{Findings and theory.}
Exact-match answer accuracy on the 1319-example GSM8K test set (greedy decoding): the $K=2$ baseline
collapses to $0.08\%$, below the base model's $47.38\%$; $K=8$ recovers to $56.03\%$ ($+8.64$ over
base) and the $K=8$ correctness-heavy reward reaches $58.53\%$ ($+11.14$ over base). Per-run seeds,
steps, and 95\% CIs are in Table~\ref{tab:controlled-comparison}; all six paired-difference CIs are in
Table~\ref{tab:paired-diffs} (Appendix~\ref{app:results}).

The reward-hacking signature is clear: format pass-rate holds across the two $K=8$ runs
($94.8\%\to94.3\%$, difference not separable from 0) while answer accuracy rises. Yet contrary to
Modification~2's intent the reward change under-delivered: its $+2.50$ gain over the matched $K=8$ run
is not separable from zero, so the binding constraint is not reward format. All checkpoints still
answer in the training format on unrelated prompts (Appendix~\ref{app:general-prompts}). The $K=2$
collapse is the within-group-variance pathology from the lectures: with $\sigma_r$ often zero or noisy,
the standardised advantage carries little signal and the policy drifts (the KL spike). Raising $K$
conditions $\sigma_r$ and raises $K_{\mathrm{eff}}$, lowering gradient-estimator variance and keeping
updates inside the KL budget.

\paragraph{Limitations.}
Each setting is a single seed, and both $K=8$ runs share phase~1 (3364 default-reward steps) stitched
onto each run's $\sim$2.5k-step phase~2; a clean reward claim needs a second seed.
